\documentclass[11pt,letterpaper]{article}
\usepackage[margin=0.45in]{geometry}\usepackage{amsmath}\usepackage{graphicx}\usepackage[final,expansion=false]{microtype}\usepackage{setspace}\usepackage[table]{xcolor}\usepackage[colorlinks=true,urlcolor=blue!60!black]{hyperref}\usepackage{charter}
\definecolor{lead}{HTML}{B85042}\definecolor{navy}{HTML}{1E2761}\definecolor{boxbg}{HTML}{EEF3FA}
\newcommand{\lead}[1]{{\sffamily\bfseries\color{lead}#1}}
\setstretch{0.97}\small\pagestyle{empty}\setlength{\parindent}{0pt}\setlength{\parskip}{3pt}
\begin{document}
\begin{center}{\Large\bfseries\color{navy}Will the AI build-out pay for itself?}\\[2pt]
{\normalsize The one-page version of \emph{The Macroeconomic Limits of AI Compute}, version 28 (31 October 2026)}\\[2pt]
{\small Desmond Fitzpatrick, Ph.D. \quad Founder, Nvara.ai \quad Retired, IBM Systems}\end{center}

\colorbox{boxbg}{\begin{minipage}{0.97\textwidth}\vspace{3pt}
\textbf{The news, and the story behind it.} On 29--30 October AI stocks lost well over \$1 trillion of value: Alphabet about \$500B on Thursday, technology stocks about \$1T on Friday (Nasdaq $-$2.9\%, chip stocks $-$7\%, Intel $-$29\%); the Nasdaq-100 lost about \$2T in October. Alphabet, Amazon and Meta reported negative free cash flow, and Meta and Alphabet are selling up to \$180B of new shares to fund the spending. As one strategist put it, the market is ``increasingly questioning whether shareholders will earn an adequate return on the enormous capital being deployed.'' That is exactly what this paper counts --- and its answer is below. Counted in today's dollars, the hole on flat spending is about \$2T, the same size as what the market marked down in October; if the builders keep raising spending, as they say they will, it is nearly twice that.
\vspace{2pt}\end{minipage}}

\lead{The question.} This year Microsoft, Alphabet, Amazon, Meta and Oracle will spend about \$780 billion on AI data centers. Musk's SpaceX/xAI, CoreWeave, Nebius and others add about \$140 billion more. Will the money come back?

\lead{How we count.} A Fermi estimate: many small factors, each bounded and sourced, with the uncertainty carried to the result. All 53 inputs are public.

\colorbox{boxbg}{\begin{minipage}{0.97\textwidth}\vspace{3pt}
\textbf{The answer in five lines}
\begin{itemize}\setlength\itemsep{1pt}
\item \textbf{The cost.} The buildings and chips cost about \$890 billion a year to own by 2030 --- wear, running costs and a return on the money --- even if spending stays flat.
\item \textbf{The revenue.} AI revenue is about \$179 billion this year. To cover the cost by 2030 it must grow 57\% a year, every year. Counted sector by sector it grows about 33\%.
\item \textbf{The gap.} About \$420 billion a year short in 2030. Flat spending is the optimistic case: every company pointed 2027 spending up.
\item \textbf{The wait.} Counting the early losses and the cost of the money, no spending path earns an ordinary return by 2035. Even a one-third cut turns a yearly profit only from 2032.
\item \textbf{What would prove it wrong.} AI revenue above \$440 billion in calendar 2028.
\end{itemize}\vspace{2pt}\end{minipage}}

\lead{New in this version.} \textbf{Every builder counted.} The rent that SpaceX/xAI, CoreWeave and the others earn from AI labs and clouds was always in the revenue we counted; their buildings and chips were not in the cost. Now they are. That one correction moves the 2030 gap from about \$260 billion (the five largest alone) to about \$420 billion. Builders that publish no figure yet (Meta's Hyperion, Crusoe, Lambda, the Gulf's Humain and Stargate UAE and others) are listed with what will bring each in; each would widen the gap.

\lead{Also new.} \emph{Inflation}, counted once: after it, flat spending is losing about 6.5\% a year in 2035. \emph{Tax}: the new write-off delays tax rather than cancelling it, so tax widens the gap. \emph{The October earnings}: Anthropic's first filed numbers matched the paper's estimate.

\lead{Why it slows.} Time saved is not money saved; checking AI's work takes time; much of the value goes to buyers and integrators; and the grid and local consent hold projects back.

\lead{What it means.} Not a crash, and not the end of AI. The likely path is write-downs and a spending retreat in 2027--29, then a cheaper, more distributed build. This week moved the timing earlier, not the destination. The sign to watch: the first builder to guide 2027 spending down.

\lead{Where I was wrong.} Earlier versions counted the other builders' rent but not their cost. It is corrected and listed, as every error will be.

{\footnotesize\sloppy\textbf{Read more.}\\
Paper, version 28: \url{https://doi.org/[v28 DOI]} \quad All versions: \url{https://doi.org/10.5281/zenodo.23406475}\\
Deck (19 slides): on the Zenodo record \quad Model, inputs and code: \url{https://github.com/desfitz/ai-compute-limits}\\
What changed: \url{https://github.com/desfitz/ai-compute-limits/blob/main/CHANGES.md} \quad First open question: \url{https://github.com/desfitz/ai-compute-limits/issues/1}\\
SSRN: \url{https://papers.ssrn.com/sol3/papers.cfm?abstract_id=7620359} \quad Comments: desmond.fitzpatrick@nvara.ai\\
Data Democracy: \url{https://datademocracy.substack.com/p/the-macroeconomic-limits-of-ai-compute}}
\end{document}
